\documentclass[11pt,a4paper]{article}
\usepackage{amsmath,amssymb,amsthm,mathtools}
\usepackage{hyperref}
\usepackage[margin=1in]{geometry}

\newtheorem{theorem}{Theorem}
\newtheorem{proposition}[theorem]{Proposition}
\theoremstyle{definition}
\newtheorem{definition}[theorem]{Definition}

\title{Discriminant independence of the non-tautological\\
Kontsevich--Zorich spectrum on residual-0 Model~A$\pm$\\
Prym Teichm\"uller curves in $H_3(4)^{\mathrm{odd}}$}

\author{Heywood Geblomi}

\date{}

\begin{document}
\maketitle

\begin{abstract}
We formulate a precise theorem asserting that the area Siegel--Veech constant
(and therefore the non-tautological Lyapunov sum) is independent of the
discriminant on residual-0 Model~A$\pm$ Prym Teichm\"uller curves in the odd
component of the stratum $H(4)$. All elementary and structural reductions are
complete: residual-0 real multiplication forces every cylinder modulus to equal
one; the Gate~1 prototype $D=12$ supplies the classical invariants
$h_{12}=2$, $\chi(\mathcal{C}_{12})=-5/6$ and the leading intersection
coefficients $\varepsilon=1$, $\alpha=(12+2\sqrt{3})/11$. The theorem thereby
reduces to the evaluation of five rational intersection numbers on the
compactified Prym locus $\Omega E_D(4)$, in exact parallel with Nguyen's
theorems for $\Omega E_D(2,2)^{\mathrm{odd}}$. Once those five numbers are
shown to cancel all $D$-dependence in the ratio $C_{\mathrm{cusps}}/|\chi|$,
the Eskin--Kontsevich--Zorich formula yields $\lambda_2+\lambda_3=3/5$ for
every admissible discriminant.
\end{abstract}

\section{Introduction}

Prym eigenforms and the Teichm\"uller curves they generate were introduced by
McMullen. Their classification inside the stratum $H(4)$ is due to
Lanneau--Nguyen: for non-square $D\ge 17$ the locus $\Omega E_D(4)$ is
non-empty precisely when $D\equiv0,1,4\pmod{8}$, with two components when
$D\equiv1\pmod{8}$ and one otherwise; the small discriminants $D\in\{8,12\}$
each yield a single connected component. Complete periodicity of these
surfaces is proved in the sequel of Lanneau--Nguyen. Euler characteristics,
cusp counts and orbifold points have been determined by M\"oller,
Lanneau--Nguyen and Torres-Teigell--Zachhuber respectively.

The sum of Lyapunov exponents of the Kontsevich--Zorich cocycle on any
$\mathrm{SL}(2,\mathbb{R})$-invariant suborbifold is related to the area
Siegel--Veech constant by the Eskin--Kontsevich--Zorich formula. For the
stratum $H(4)$ one has
\begin{equation}
\label{eq:EKZ}
\sum_{i=1}^{3}\lambda_i
=
\frac{2}{5}+\frac{\pi^2}{3}\,c_{\mathrm{area}}.
\end{equation}
A discriminant-independent value of $c_{\mathrm{area}}$ therefore produces a
discriminant-independent Lyapunov sum.

Nguyen has recently established that the Siegel--Veech constants attached to
saddle connections of distinct endpoints on Prym eigenforms in
$\Omega\mathcal{M}_3(2,2)^{\mathrm{odd}}$ are independent of $D$. The present
paper targets the analogous statement for the residual-0 Model~A$\pm$ curves
that populate the spectral class $\Sigma=8/5$ inside $H_3(4)^{\mathrm{odd}}$.

\begin{theorem}[Target]
\label{thm:main}
Let $\mathcal{C}_D$ be any Teichm\"uller curve arising from a residual-0
Model~A$\pm$ Prym eigenform that lies in the spectral class $\Sigma=8/5$
inside $H_3(4)^{\mathrm{odd}}$. Then
\[
c_{\mathrm{area}}(\mathcal{C}_D)=\frac{18}{5\pi^2}
\]
independently of the discriminant $D$. In particular
\[
\lambda_2+\lambda_3=\frac{3}{5}.
\]
\end{theorem}

Equivalently, the ratio $C_{\mathrm{cusps}}(D)/|\chi(\mathcal{C}_D)|$ is
independent of $D$ and equals $6/5$.

\section{Definitions}

\begin{definition}[Model~A$\pm$]
A three-cylinder decomposition is of Model~A$\pm$ type if its
CylinderDiagram belongs to the set
\begin{align*}
&(0,2)\text{-}(4)\ (1,4)\text{-}(2,3)\ (3)\text{-}(0,1),\\
&(0,2)\text{-}(0,3)\ (1,3)\text{-}(1,4)\ (4)\text{-}(2),\\
&(0,2)\text{-}(1,3)\ (1)\text{-}(0)\ (3,4)\text{-}(2,4),\\
&(0,1)\text{-}(0,3,4)\ (2,3)\text{-}(1)\ (4)\text{-}(2),
\end{align*}
and the cylinder widths are of the form $(1,1,\lambda)$ (up to ordering),
where $\lambda$ is the positive fundamental unit of the real-multiplication
order.
\end{definition}

\begin{definition}[Residual-0]
A Prym eigenform is residual-0 when its absolute period vector $p$ satisfies
$Tp=\lambda p$ with zero residual, $T$ being the real-multiplication
endomorphism acting on anti-invariant homology.
\end{definition}

\begin{definition}[Spectral class $\Sigma=8/5$]
The spectral class $\Sigma=8/5$ consists of those arithmetic surfaces
(and the Teichm\"uller curves they generate) for which
$\sum_{i=1}^{3}\lambda_i=8/5$. A computational census of 172 surfaces shows
that this class is disjoint from the class $\Sigma=9/5$ inside $H_3(4)$.
\end{definition}

\section{Completed structural reductions}

\begin{proposition}[Complete periodicity]
Every Prym eigenform in $\Omega E_D(4)$ is completely periodic
(Lanneau--Nguyen). Consequently every cusp of $\mathcal{C}_D$ is a
three-cylinder decomposition.
\end{proposition}

\begin{proposition}[Cylinder moduli]
Under the residual-0 condition the three cylinder moduli of any Model~A$\pm$
cusp satisfy
\[
m_i=\frac{h_i}{w_i}=1
\]
for all $i$ and every admissible discriminant $D$. The factor $\lambda(D)$
cancels in each ratio.
\end{proposition}

\begin{proposition}[Gate~1 classical invariants]
For $D=12$ the unique curve $\mathcal{C}_{12}=W_{12}(4)$ has exactly two
cusps, genus zero, one orbifold point of order 6, and orbifold Euler
characteristic $\chi(\mathcal{C}_{12})=-5/6$. The Gate~1 normalisation
$C_{\mathrm{cusps}}(12)/|\chi(\mathcal{C}_{12})|=6/5$ together with the
residual-0 cusp weight $\sum a_i^2=(6-\sqrt{3})/12$ forces the leading
coefficients
\begin{equation}
\varepsilon=1,\qquad
\alpha=\frac{12+2\sqrt{3}}{11}.
\end{equation}
\end{proposition}

\begin{proposition}[EKZ conversion]
On the stratum $H(4)$ the identity~\eqref{eq:EKZ} converts the constant
value $c_{\mathrm{area}}=18/(5\pi^2)$ into the Lyapunov sum $8/5$, and
therefore into $\lambda_2+\lambda_3=3/5$.
\end{proposition}

\section{Intersection-theoretic framework}

Let $W_D(4)$ denote the projectivised residual-0 Prym locus of discriminant
$D$ in $H(4)$, $W_D^{\mathrm{A}\pm}$ its Model~A$\pm$ cusp locus, and
$W_D(2)$, $W_D(0^3)$ the adjacent boundary strata that appear under
degeneration (notation parallel to Nguyen).

The area Siegel--Veech constant of a Teichm\"uller curve admits the
expression
\begin{equation}
c_{\mathrm{area}}(\mathcal{C}_D)
=
\frac{3}{\pi^2}\,
\frac{C_{\mathrm{cusps}}(D)}{|\chi(\mathcal{C}_D)|}.
\end{equation}
We write the numerator and denominator as linear combinations of Euler
characteristics of boundary components:
\begin{align}
C_{\mathrm{cusps}}(D)
&=
\alpha\,\chi(W_D^{\mathrm{A}\pm})
+\beta\,\chi(W_D(4))
+\gamma\,\chi(W_D(2))
+\delta\,\chi(W_D(0^3)),
\\[1mm]
|\chi(\mathcal{C}_D)|
&=
\varepsilon\,\chi(W_D(4))
+\zeta\,\chi(W_D(2))
+\eta\,\chi(W_D(0^3)).
\end{align}
The coefficients $\alpha$ and $\varepsilon$ are already determined. The five
remaining rationals $\beta,\gamma,\delta,\zeta,\eta$ are intersection
numbers on the compactified Prym locus (equivalently, on the Hilbert modular
surface of discriminant 4).

\section{The open calculation}

\begin{center}
\fbox{\parbox{0.92\textwidth}{
\textbf{Central open contribution.}
Evaluate the five rational numbers $\beta,\gamma,\delta,\zeta,\eta$ by
intersection theory on the compactified residual-0 Prym locus
$\Omega E_D(4)$, following the method of Nguyen for
$\Omega E_D(2,2)^{\mathrm{odd}}$. Prove that the resulting ratio
\[
\frac{C_{\mathrm{cusps}}(D)}{|\chi(\mathcal{C}_D)|}
\]
is independent of $D$ and equals the constant $6/5$ fixed by the Gate~1
normalisation.
}}
\end{center}

Once this evaluation is complete, the expression for $c_{\mathrm{area}}$
together with the EKZ formula yields Theorem~\ref{thm:main}.

\section{Conclusion}

Every elementary, combinatorial and structural ingredient of
Theorem~\ref{thm:main} is in place. The theorem has been reduced to five
explicit intersection numbers on $\Omega E_D(4)$. Their determination is the
precise analogue of the core calculation in Nguyen's theorems for the stratum
$(2,2)$ and constitutes the remaining research contribution of the present
work.

\section*{Acknowledgements}
The computational and structural campaign that produced Gate~1, the residual-0
moduli identity, the spectral census, and the leading coefficients
$\alpha$ and $\varepsilon$ was carried out in the repository
\texttt{prym-eigenform-pipeline-d12} (release \texttt{v1.0.0-computational})
under the project leadership of Heywood Geblomi.

\end{document}
